
%==================================================
%%%%%%%%%%%%%%%%%%%%%%%%%%%%%%%%%%%%%%%%%%%%%%%%%%%%
\chapter{Gestion du projet : Analyse et Planification}

\section{Introduction}
Nous allons, dans ce chapitre, approfondir les besoins qui motivent le développement du système Smart Focus. C’est une étape clé pour s’assurer que la solution finale réponde aux attentes des utilisateurs et aux contraintes techniques identifiées.

En faisant partie de l’équipe projet et en assurant les différentes missions qui nous ont été assignées, nous avons pris une part active à cette phase d’analyse. Nous verrons tout d’abord comment les besoins ont été collectés, puis analysés et priorisés. Nous présenterons ensuite les acteurs du système, les diagrammes de cas d’utilisation ainsi que l’architecture globale adoptée. Enfin, nous rentrerons dans le détail de la planification des sprints, afin de poser les bases du projet en alignant clairement les objectifs fonctionnels et techniques sur les attentes réelles des utilisateurs.
%%%%%%%%%%%%%%%%%%%%%%%%%%%%%%%%%%%%%%%%%%%%%%%%%%%%
% RECUEIL DES BESOINS
%%%%%%%%%%%%%%%%%%%%%%%%%%%%%%%%%%%%%%%%%%%%%%%%%%%%
\section{Recueil des besoins}

La collecte des besoins constitue une étape essentielle dans la conception du système \textit{SmartFocus}. Elle permet de comprendre les attentes des utilisateurs, d’identifier les fonctionnalités nécessaires et de définir les contraintes du projet.

Cette démarche vise à poser les bases d’une suite solide du projet, en alignant clairement les objectifs fonctionnels et techniques avec les attentes réelles. Pour obtenir une vision claire et réaliste, plusieurs approches complémentaires ont été adoptées.
% IDENTIFICATION DES ACTEURS
%%%%%%%%%%%%%%%%%%%%%%%%%%%%%%%%%%%%%%%%%%%%%%%%%%%%
\subsection{Identification des acteurs}
Le système \textit{Smart Focus} met en jeu plusieurs acteurs, qu’ils soient principaux ou secondaires. Le tableau suivant présente de manière synthétique le rôle de chacun au sein du système.


\begin{table}[H]
\centering
\renewcommand{\arraystretch}{1.5}
\setlength{\tabcolsep}{8pt}

\begin{tabularx}{\textwidth}{|>{\raggedright\arraybackslash}p{4cm}|c|>{\raggedright\arraybackslash}X|}
\hline
\textbf{Acteur} & \textbf{Type} & \textbf{Description} \\
\hline

Utilisateur 
& Principal 
& Étudiant, professionnel ou enseignant utilisant le système pour suivre sa concentration, gérer ses sessions et consulter les informations fournies. \\
\hline

Service d’intelligence artificielle externe 
& Secondaire 
& Service externe accessible via API, utilisé pour analyser les données, générer des recommandations et fournir des fonctionnalités intelligentes (chatbot, assistance). \\
\hline

\end{tabularx}

\caption{Identification des acteurs du système}
\label{tab:acteurs}

\end{table}
%\subsection{Identification des besoins}

%En se basant sur les méthodes de recueil décrites ci-dessus, il a été possible d’identifier plusieurs besoins clés pour le projet \textit{Smart Focus}. Ces besoins traduisent à la fois les attentes des utilisateurs en matière de fonctionnalités et les contraintes techniques à respecter.

%\vspace{0.5em}
%\begin{table}[H]
 %   \centering
  %\begin{tabular}{|p{5.5cm}|p{8.5cm}|}
   %     \hline
      %  \textbf{Catégorie} & \textbf{Description} \\ \hline
      %  Détection de la concentration & Analyse en temps réel de la posture, de la fatigue et des distractions via la vision par ordinateur pour calculer un score de concentration. \\ \hline
       % Planification intelligente & Génération automatique de sessions d'étude adaptées au profil de l'utilisateur, à son sommeil et à ses examens à venir, en exploitant l'\ac{IA}. \\ \hline
      %  Assistance pédagogique & Chatbot \ac{RAG} permettant de poser des questions sur ses cours (PDF), de générer des quiz QCM et des flashcards avec répétition espacée. \\ \hline
      %  Suivi du sommeil & Enregistrement et analyse de la qualité du sommeil, avec adaptation automatique de la charge de travail du lendemain. \\ \hline
     %   Feedback intelligent & Retours visuels (LEDs, écran TFT), sonores et via notifications mobiles pour alerter l'utilisateur en cas de baisse de concentration ou de mauvaise posture. \\ \hline
     %   Interface mobile ergonomique & Application Flutter intuitive offrant dashboard, statistiques, gestion des sessions et accès à toutes les fonctionnalités du système. \\ \hline
 %   \end{tabular}
 %   \caption{Synthèse des besoins identifiés}
 %   \label{tab:besoins_identifies}
%\end{table}


%%%%%%%%%%%%%%%%%%%%%%%%%%%%%%%%%%%%%%%%%%%%%%%%%%%%
% ANALYSE DES BESOINS
%%%%%%%%%%%%%%%%%%%%%%%%%%%%%%%%%%%%%%%%%%%%%%%%%%%%
\subsection{Identification des besoins}

Pour mieux structurer les exigences du projet \textit{Smart Focus}, les besoins identifiés ont été classés selon leur nature et leur rôle dans le fonctionnement global du système. Cette classification facilite la compréhension des priorités et des contraintes à prendre en compte.

\vspace{0.5em}
\begin{itemize}
    \item \textbf{Besoins fonctionnels} :
    \begin{itemize}
        \item Capturer et analyser les données visuelles afin d'évaluer la posture, la fatigue et le niveau de concentration de l'utilisateur en temps réel.
        \item Détecter les situations de distraction et générer des alertes adaptées.
        \item Calculer un score de concentration dynamique basé sur plusieurs indicateurs comportementaux.
        \item Gérer des sessions de travail intelligentes, incluant le démarrage, le suivi et l'historique.
        \item Proposer un système de planification adaptatif permettant d'organiser les sessions en fonction de l'état de l'utilisateur (fatigue, concentration, sommeil) et de ses examens.
        \item Intégrer un assistant pédagogique basé sur un chatbot \ac{RAG} permettant d'exploiter des documents (PDF et CSV), de répondre aux questions, de générer des quiz et des flashcards.
        \item Offrir un système de suivi du sommeil permettant d'évaluer sa qualité et d'influencer l'organisation des activités.
        \item Fournir un accompagnement via un agent vocal intelligent capable d'interagir avec l'utilisateur et de proposer des recommandations.
        \item Afficher les données et statistiques à travers une application mobile.
        \item Assurer la synchronisation des données entre les différents composants du système (boîtier \ac{IoT}, backend, application).
    \end{itemize}

    \item \textbf{Besoins non fonctionnels} :
    \begin{itemize}
        \item \textbf{Performance :} Traitement rapide des données garantissant un fonctionnement en temps réel.
        \item \textbf{Précision :} Analyses IA fiables et cohérentes.
        \item \textbf{Ergonomie :} Interface simple, intuitive et adaptée à une utilisation quotidienne.
        \item \textbf{Disponibilité :} Fonctionnement maintenu même en cas de perte temporaire de connexion.
        \item \textbf{Scalabilité :} Architecture permettant l'ajout de nouvelles fonctionnalités.
        \item \textbf{Sécurité :} Protection des données de l'utilisateur (\ac{JWT}, bcrypt).
        \item \textbf{Compatibilité :} Fonctionnement sur plateforme embarquée et accès via application mobile.
        \item \textbf{Maintenabilité :} Code structuré facilitant les mises à jour et améliorations.
    \end{itemize}

    \item \textbf{Contraintes} :
    \begin{itemize}
        \item Ressources matérielles limitées de Raspberry (mémoire, puissance de calcul).
        \item Dépendance à un service \ac{LLM} externe (Groq, avec Google Gemini en alternative) nécessitant une connexion Internet.
        \item Séparation des responsabilités entre deux membres de l'équipe (hardware / logiciel).
        \item Compatibilité avec les technologies retenues (FastAPI, Flutter, PostgreSQL, ChromaDB).
    \end{itemize}
\end{itemize}
\vspace{1em}

\noindent\colorbox{gray!15}{%
    \parbox{\dimexpr\linewidth-2\fboxsep}{%
        \small
        \textbf{Résumé :} Cette classification permet de mieux cibler les efforts de développement en mettant en lumière les aspects essentiels à couvrir pour assurer la réussite du projet.
    }%
}

\subsection{Priorisation des besoins}

Afin d'assurer le progrès efficace du projet \textit{Smart Focus}, il est primordial de classer les besoins selon leur influence sur les utilisateurs et leur viabilité technique. Cette phase permet de focaliser les efforts sur les fonctionnalités qui sont les plus cruciales.

\vspace{0.5em}
\begin{table}[H]
    \centering
    \renewcommand{\arraystretch}{1.3}
    
    \begin{tabular}{|p{5cm}|p{6.5cm}|c|}
        \hline
        \textbf{Besoin} & \textbf{Justification} & \textbf{Priorité} \\ \hline
        
        Authentification et gestion de profil 
        & Nécessaire pour la gestion des utilisateurs et la persistance des données 
        & Haute \\ \hline
        
        Détection de la concentration (score focus) 
        & Cœur du système basé sur la vision par ordinateur et le hardware 
        & Haute \\ \hline
        
        Chatbot \ac{RAG} (upload, chat, quiz, flashcards) 
        & Apporte une forte valeur ajoutée pédagogique et personnalisation via l’IA 
        & Haute \\ \hline
        
        Planification intelligente 
        & Permet d’adapter les sessions en fonction des données utilisateur et du niveau de fatigue 
        & Haute \\ \hline
        
        Dashboard et statistiques 
        & Permet l’analyse des performances et l’amélioration continue de l’utilisateur 
        & Moyenne \\ \hline
        
        Posture et feedback ergonomique 
        & Améliore l’expérience utilisateur via des recommandations physiques basées sur la vision 
        & Moyenne \\ \hline
        
        Suivi du sommeil et alarme 
        & Fonction complémentaire influençant indirectement la performance de l’utilisateur 
        & Basse \\ \hline
        
        Agent vocal intelligent
        & Fonctionnalité complémentaire prévue en phase finale, permettant une interaction naturelle avec le système
        & Basse \\ \hline
        
    \end{tabular}

    \caption{Priorisation des besoins}
    \label{tab:priorisation_besoins}
\end{table}

\vspace{1em}
\noindent\colorbox{gray!15}{%
    \parbox{\dimexpr\linewidth-2\fboxsep}{%
        \textbf{Conclusion :} On a accordé la priorité aux besoins relatifs à l’authentification, à l’assistance pédagogique et à la planification intelligente, afin de poser une base solide pour le système tout en prévoyant une intégration progressive des fonctionnalités dépendant du hardware.
    }%
}

%%%%%%%%%%%%%%%%%%%%%%%%%%%%%%%%%%%%%%%%%%%%%%%%%%%%
% MODÉLISATION — CAS D'UTILISATION
%%%%%%%%%%%%%%%%%%%%%%%%%%%%%%%%%%%%%%%%%%%%%%%%%%%%
\section{Modélisation des besoins}

\subsection{Diagramme de cas d'utilisation global}
Le diagramme suivant, présente la liste de tous les cas d'utilisation du système Smart Focus, classés par module fonctionnel. Le système comporte 46 cas d’usage organisés en 9 modules.
\begin{figure}[htbp]
    \centering
    \includegraphics[width=1\textwidth]{images/UCDiagram.png}
    \caption{Diagramme de cas d'utilisation global}
    \label{fig:uml_global}
\end{figure}


\section{Planification du projet}
\subsection{Backlog du produit}
{\singlespacing
\begin{table}[H]
\centering
\small
\renewcommand{\arraystretch}{1.2}
\setlength{\tabcolsep}{5pt}

\begin{tabularx}{\textwidth}{|c|p{3cm}|X|c|}
\hline
\textbf{ID} & \textbf{Fonctionnalité} & \textbf{User Story} & \textbf{Priorité} \\
\hline

US1 & Authentification
& En tant qu’utilisateur, je veux créer un compte et me connecter afin d’accéder à mes données personnelles.
& Haute \\ \hline

US2 & Démarrage session
& En tant qu’utilisateur, je veux démarrer une session de travail afin de suivre mon niveau de concentration.
& Haute \\ \hline

US3 & Capture vidéo (Raspberry Pi)
& En tant que système, je veux capturer les images en temps réel afin d’analyser le comportement de l’utilisateur.
& Haute \\ \hline

US4 & Détection de concentration
& En tant que système, je veux analyser les images pour calculer un score de concentration.
& Haute \\ \hline

US5 & Détection fatigue et posture
& En tant que système, je veux détecter la fatigue et la posture afin d’évaluer l’état de l’utilisateur.
& Haute \\ \hline

US6 & Enregistrement des métriques
& En tant que système, je veux stocker les données de session afin de permettre leur exploitation.
& Haute \\ \hline

US7 & Génération d’alertes
& En tant que système, je veux déclencher des alertes lorsque la fatigue est élevée afin d’améliorer la concentration.
& Haute \\ \hline

US8 & Affichage temps réel
& En tant qu’utilisateur, je veux visualiser mes scores en temps réel sur l’application mobile.
& Haute \\ \hline

US9 & Chatbot intelligent (RAG)
& En tant qu’utilisateur, je veux interagir avec un chatbot pour poser des questions et réviser mes cours.
& Haute \\ \hline

US10 & Import de documents
& En tant qu’utilisateur, je veux importer mes cours afin qu’ils soient utilisés par le chatbot.
& Haute \\ \hline

US11 & Génération de quiz/flashcards
& En tant qu’utilisateur, je veux générer des quiz pour tester mes connaissances.
& Moyenne \\ \hline

US12 & Planification intelligente
& En tant qu’utilisateur, je veux obtenir un planning adapté à mon niveau de fatigue et mes performances.
& Haute \\ \hline

US13 & Dashboard et statistiques
& En tant qu’utilisateur, je veux visualiser mes performances pour suivre mon évolution.
& Moyenne \\ \hline

US14 & Feedback ergonomique
& En tant qu’utilisateur, je veux recevoir des recommandations sur ma posture.
& Moyenne \\ \hline

US15 & Suivi du sommeil
& En tant qu’utilisateur, je veux suivre mon sommeil afin d’améliorer mes sessions de travail.
& Basse \\ \hline

US16 & Agent vocal
& En tant qu’utilisateur, je veux interagir avec le système via la voix.
& Basse \\ \hline

\end{tabularx}

\caption{Backlog du produit}
\label{tab:backlog}
\end{table}
}
%%%%%%%%%%%%%%%%%%%%%%%%%%%%%%%%%%%%%%%%%%%%%%%%%%%%%%%%%%%%%%%%%%%%%%%%%%%%%%%%%%%%%%%%%%%%%%%%%%%%%%%%%%%%%%%%%%%%%%%%%%%%%%%%%%


\subsection{Planification des sprints}

Le développement du projet a été organisé en quatre sprints successifs, chacun couvrant une période d’environ un mois. Cette planification a été définie de manière progressive afin d’assurer un équilibre entre les différentes composantes du système, notamment l’application mobile, les modules intelligents et la vision par ordinateur.
{\singlespacing
\begin{table}[H]
\centering
\small
\renewcommand{\arraystretch}{1.2}
\setlength{\tabcolsep}{5pt}

\begin{tabularx}{\textwidth}{|c|p{3.8cm}|X|c|}
\hline
\textbf{Sprint} & \textbf{Intitulé} & \textbf{Description} & \textbf{Durée} \\
\hline

Sprint 1 & 
Authentification et interfaces de base & 
Développement du module d’authentification permettant l’inscription, la connexion et la gestion du profil utilisateur côté backend et application mobile. Mise en place des interfaces principales de l’application (tableau de bord, statistiques, sommeil, paramètres) ainsi que de la navigation entre les écrans. &
20 jours \\ \hline

Sprint 2 & 
Surveillance intelligente de la concentration (Vision + Score) & 
Implémentation du module de vision par ordinateur permettant l’analyse en temps réel de l’état de l’utilisateur à partir du flux vidéo. Détection de la fatigue, des distractions et de la posture via des techniques d’intelligence artificielle, puis calcul d’un score global de concentration transmis au backend. &
40 jours \\ \hline

Sprint 3 & 
Planning intelligent et assistant d’apprentissage & 
Développement du système de planning intelligent basé sur l’intelligence artificielle pour la génération automatique des sessions d’étude. Intégration du chatbot RAG permettant l’interrogation de documents PDF, ainsi que la génération automatique de quiz et de flashcards avec révision espacée (SM-2). Connexion des interfaces mobiles aux données dynamiques du backend. &
30 jours \\ \hline

Sprint 4 & 
Système IoT et alertes intelligentes & 
Intégration du dispositif matériel de monitoring permettant la capture continue des données utilisateur via caméra et capteurs. Mise en place du suivi temps réel dans l’application mobile ainsi que du système d’alertes intelligentes basé sur le niveau de concentration et les données de sommeil afin d’adapter le comportement du système et le planning de travail. &
30 jours \\ \hline

\end{tabularx}

\caption{Planification des sprints du projet}
\label{tab:sprints_final}
\end{table}
}
%==================================================
\section{Diagramme de Gantt}

\begin{figure}[H]
\centering
\includegraphics[width=\textwidth]{images/Gant.png}
\caption{Diagramme de Gantt}
\end{figure}
%==================================================
\section{Environnement de travail}

\subsection{Choix des technologies}

Le développement du système \textit{Smart Focus} repose sur un ensemble de technologies modernes permettant de garantir performance, modularité et évolutivité. Les technologies retenues sont présentées ci-dessous, regroupées par catégorie.

\subsubsection{Langages de programmation}

\begin{itemize}

\item \textbf{Python} \cite{python} \textbf{:} Langage de programmation largement utilisé dans le développement web, l'intelligence artificielle et le traitement d'images. Dans notre projet, Python est utilisé pour développer le backend avec FastAPI, le pipeline de vision par ordinateur pour l'analyse en temps réel via caméra, ainsi que les modules d'intelligence artificielle.

\begin{figure}[H]
    \centering
    \includegraphics[width=0.15\textwidth]{images/pythonLogo.png}
    \caption{Logo Python}
\end{figure}

\item \textbf{Dart} \cite{dart} \textbf{:} Langage de programmation développé par Google, optimisé pour la création d'interfaces utilisateur performantes et multiplateformes. Dans notre projet, Dart est le langage utilisé avec le framework Flutter pour développer l'application mobile Smart Focus.

\begin{figure}[H]
    \centering
    \includegraphics[width=0.3\textwidth]{images/dart.jpg}
    \caption{Logo Dart}
\end{figure}

\end{itemize}

\subsubsection{Frameworks et bibliothèques}

\begin{itemize}

\item \textbf{FastAPI} \cite{fastapi} \textbf{:} Framework Python moderne conçu pour la création d'API REST performantes avec documentation automatique. Dans notre projet, FastAPI gère l'ensemble des services backend : authentification, gestion des sessions de monitoring, planning intelligent, chatbot, quiz, flashcards et suivi du sommeil.

\begin{figure}[H]
    \centering
    \includegraphics[width=0.3\textwidth]{images/fastapi.jpg}
    \caption{Logo FastAPI}
\end{figure}

\item \textbf{Flutter} \cite{flutter} \textbf{:} Framework open-source de Google pour le développement d'applications mobiles multiplateformes à partir d'un seul code source. Dans notre projet, Flutter est utilisé pour développer l'application mobile qui permet à l'utilisateur de consulter son tableau de bord, gérer son planning, interagir avec le chatbot, réviser ses flashcards et suivre ses sessions de monitoring en temps réel.

\begin{figure}[H]
    \centering
    \includegraphics[width=0.2\textwidth]{images/flutter.png}
    \caption{Logo Flutter}
\end{figure}



\item \textbf{OpenCV} \cite{opencv} \textbf{:} Bibliothèque open-source de traitement d'images et de vision par ordinateur. Dans notre projet, OpenCV est utilisé au sein du module pi\_client pour l'acquisition du flux vidéo depuis la caméra et le prétraitement des images avant analyse.

\begin{figure}[H]
    \centering
    \includegraphics[width=0.2\textwidth]{images/openCV.png}
    \caption{Logo OpenCV}
\end{figure}

\item \textbf{MediaPipe} \cite{mediapipe} \textbf{:} Framework de Google pour l'analyse en temps réel de flux multimédia. Dans notre projet, MediaPipe est utilisé dans le pipeline de vision pour la détection de la pose corporelle et des repères faciaux, permettant le calcul de quatre scores comportementaux : attention, posture, vigilance et focus global.

\begin{figure}[H]
    \centering
    \includegraphics[width=0.45\textwidth]{images/mediaPipe.png}
    \caption{Logo MediaPipe}
\end{figure}

\item \textbf{LangChain} \cite{langchain} \textbf{:} Framework Python pour le développement d'applications basées sur des modèles de langage. Dans notre projet, LangChain orchestre le pipeline RAG du chatbot : découpage des documents PDF en chunks, génération des embeddings, recherche sémantique et construction du prompt contextuel envoyé au LLM.

\begin{figure}[H]
    \centering
    \includegraphics[width=0.8\textwidth]{images/langchainLogo.png}
    \caption{Logo LangChain}
\end{figure}

\end{itemize}

\subsubsection{Intelligence artificielle et bases de données}

\begin{itemize}

\item \textbf{Groq} \cite{groq} \textbf{:} Plateforme d'inférence LLM haute performance. Dans notre projet, Groq est utilisé pour l'exécution des modèles de langage à travers l'approche \ac{RAG}, permettant au chatbot de répondre aux questions contextualisées à partir des documents indexés, de générer des quiz et des flashcards, et de personnaliser les sujets de révision dans le planning intelligent.

\begin{figure}[H]
    \centering
    \includegraphics[width=0.2\textwidth]{images/groq.png}
    \caption{Logo Groq}
\end{figure}

\item \textbf{ChromaDB} \cite{chromadb} \textbf{:} Base de données vectorielle open-source. Dans notre projet, ChromaDB stocke les embeddings des chunks de documents PDF indexés et permet la recherche par similarité cosinus lors des requêtes du chatbot, de la génération de quiz et de flashcards.

\begin{figure}[H]
    \centering
    \includegraphics[width=0.25\textwidth]{images/chromaDB.png}
    \caption{Logo ChromaDB}
\end{figure}

\item \textbf{PostgreSQL} \cite{postgresql} \textbf{:} Système de gestion de base de données relationnelle robuste et open-source. Dans notre projet, PostgreSQL stocke l'ensemble des données structurées : comptes utilisateurs, profils, sessions de travail, snapshots de monitoring, événements de concentration, sessions de planning, examens, documents, messages du chatbot, quiz, flashcards et enregistrements de sommeil.

\begin{figure}[H]
    \centering
    \includegraphics[width=0.3\textwidth]{images/PostgreSQL.jpg}
    \caption{Logo PostgreSQL}
\end{figure}

\end{itemize}

%==================================================

\subsection{Choix des logiciels}

Dans le cadre du développement du projet \textit{Smart Focus}, plusieurs outils logiciels ont été utilisés afin d'assurer un environnement de travail efficace et structuré.

\begin{itemize}

\item \textbf{Visual Studio Code} \cite{vscode} \textbf{:} Éditeur de code source léger et extensible développé par Microsoft. Dans notre projet, VS Code est l'environnement de développement principal utilisé pour l'implémentation du backend Python, de l'application mobile Flutter et du module de vision par ordinateur.

\begin{figure}[H]
    \centering
    \includegraphics[width=0.15\textwidth]{images/vs.png}
    \caption{Visual Studio Code}
\end{figure}

\item \textbf{Android Studio} \cite{androidstudio} \textbf{:} Environnement de développement intégré officiel pour Android. Dans notre projet, Android Studio est utilisé pour la gestion de l'émulateur Android lors des tests de l'application mobile Flutter, ainsi que pour la configuration des SDK nécessaires au build.

\begin{figure}[H]
    \centering
    \includegraphics[width=0.15\textwidth]{images/androidStudio.png}
    \caption{Android Studio}
\end{figure}

\item \textbf{Git et GitHub} \cite{github} \textbf{:} Système de contrôle de version distribué et plateforme d'hébergement de code. Dans notre projet, Git et GitHub assurent le versionnement du code source, avec un historique complet des modifications pour chaque composant.

\begin{figure}[H]
    \centering
    \includegraphics[width=0.25\textwidth]{images/gitHub.png}
    \caption{GitHub}
\end{figure}

\item \textbf{Docker} \cite{docker} \textbf{:} Plateforme de conteneurisation permettant d'isoler et de déployer des applications. Dans notre projet, Docker est utilisé pour conteneuriser les services du backend (FastAPI, PostgreSQL, ChromaDB), garantissant la portabilité de l'environnement et simplifiant le déploiement.

\begin{figure}[H]
    \centering
    \includegraphics[width=0.25\textwidth]{images/docker.png}
    \caption{Docker}
\end{figure}


\item \textbf{Draw.io} \cite{drawio} \textbf{:} Outil de diagrammes en ligne gratuit et open-source. Dans notre projet, Draw.io a été utilisé pour la conception des diagrammes UML : diagrammes de cas d'utilisation, de classes, de séquence et d'architecture système.

\begin{figure}[H]
    \centering
    \includegraphics[width=0.25\textwidth]{images/draw io.png}
    \caption{Draw.io}
\end{figure}

\item \textbf{Overleaf} \cite{overleaf} \textbf{:} Plateforme \LaTeX{} collaborative en ligne. Dans notre projet, Overleaf est utilisé pour la rédaction du rapport de projet, permettant un formatage professionnel et une collaboration en temps réel sur le document.

\begin{figure}[H]
    \centering
    \includegraphics[width=0.2\textwidth]{images/overleaf.png}
    \caption{Overleaf}
\end{figure}

\end{itemize}

%==================================================
%==================================================
\section{Architecture du système}

Le système \textit{Smart Focus} repose sur une architecture multicouche visant à séparer les différentes responsabilités et à assurer une meilleure modularité. Cette organisation permet de faciliter la maintenance, l’évolution du système ainsi que l’intégration de nouvelles fonctionnalités.

L’architecture proposée est composée de quatre couches principales :

\begin{enumerate}
    \item \textbf{Couche Hardware} — Boîtier Raspberry Pi (caméra OV2640, capteurs, écran TFT, LEDs WS2812B)

    \item \textbf{Couche Backend} — Serveur FastAPI (Python 3.11) assurant la logique métier et les services \ac{API} REST

    \item \textbf{Couche Données} — PostgreSQL (données relationnelles) et ChromaDB (données vectorielles), avec stockage des fichiers (PDF)

    \item \textbf{Couche Client} — Application mobile Flutter utilisant Riverpod et GoRouter pour la gestion d’état et la navigation
\end{enumerate}




Cette architecture permet d’assurer une séparation claire entre les différentes couches du système, tout en garantissant une communication fluide entre les composants via des interfaces bien définies.
%==================================================
\section{Diagramme de classes global}

Le diagramme de classes suivant illustre la structure globale du système, en mettant en évidence les principales entités, leurs attributs ainsi que leurs relations.

\begin{figure}[H]
\centering
\includegraphics[width=1.1\textwidth]{images/DiagClassGlobale.png}
\caption{Diagramme de classes global}
\label{fig:class_diagram}
\end{figure}
%==================================================
\section*{Conclusion}

Dans ce chapitre, nous avons présenté l'analyse et la modélisation des besoins du système Smart Focus. Nous avons d'abord décrit les méthodes de recueil des besoins, puis classé et priorisé ces besoins selon leur nature et leur impact.

Nous avons ensuite identifié les acteurs du système et présenté les 46 cas d'utilisation répartis en 9 modules fonctionnels, illustrés dans le diagramme de cas d'utilisation global. L'architecture multicouche a été décrite, couvrant le boîtier, le backend FastAPI, les bases de données et l'application Flutter.

Enfin, la planification des sprints a permis de structurer le développement de manière itérative, en alignant les efforts avec les priorités identifiées.
